\chapter{爆破：定义与初步性质}\label{chap:ag-2}
\section{射影态射与丰沛层的回顾}\label{sec:ag-2-1}
% 原PDF p15--p17上部.
\begin{numpara}[\texorpdfstring{$\operatorname{Proj}$}{Proj} 的回顾]\label{item:ag-2-1-1}
设 $S=S_\bullet=\bigoplus_{n\in\mathbb N}S_n$ 是分次环. 记 $\operatorname{Proj}(S)$ 为不包含理想 $S_+:=\bigoplus_{n>0}S_n$ 的齐次素理想全体.
对 $S_\bullet$ 的齐次理想 $I$，记 $V(I):=\{\mathfrak p\in\operatorname{Proj}(S_\bullet),\ \mathfrak p\supset I\}$. 这些 $V(I)$ 是某个拓扑的闭集，称该拓扑为 Zariski 拓扑.
若 $f$ 是 $S_+$ 的齐次元素，记 $D(f):=\{\mathfrak p\in\operatorname{Proj}(S_\bullet),\ f\notin\mathfrak p\}$. 它是 $\operatorname{Proj}(S_\bullet)$ 的开集，这种形式的开集构成 Zariski 拓扑的邻域基. 对另一个齐次元素 $f'$，有 $D(ff')=D(f)\cap D(f')$. 族 $(D(f_i))_{i\in I}$ 覆盖 $\operatorname{Proj}(S_\bullet)$，当且仅当 $f_i$ 生成的理想包含 $S_+$ 中每个元素的某个幂.
设 $f$ 的次数为 $d$. 局部化环 $S_f=S_{f,\bullet}=S_\bullet[X]/(Xf-1)$ 带有由 $S_\bullet$ 的分次诱导的分次，其中 $X$ 的次数为 $-d$. 记 $S_{(f)}=S_{f,0}$ 为次数 $0$ 的元素组成的子环. 有时用归纳极限来理解更方便：有典范的 $S_\bullet$-模同构（参见 EGA 0.1.6）
\begin{align*}
S_{f,\bullet}=\varinjlim\bigl(S_\bullet\xrightarrow{\times f}S_\bullet\xrightarrow{\times f}\cdots\xrightarrow{\times f}S_\bullet\xrightarrow{\times f}\cdots\bigr).
\end{align*}
若将归纳极限看作 $S_\bullet^{(\mathbb N)}=S_\bullet[X]$ 的商，则它是 $S_\bullet$-代数同构. 通过限制，得到 $S_0$-模同构
\begin{align*}
S_{(f)}=S_{f,0}=\varinjlim\bigl(S_0\xrightarrow{\times f}S_d\xrightarrow{\times f}\cdots\xrightarrow{\times f}S_{dn}\xrightarrow{\times f}\cdots\bigr).
\end{align*}
若将归纳极限看作环 $\bigoplus_nS_{dn}$ 的商，则它是 $S_0$-代数同构. 当 $d=1$ 时，后一表达式给出同构 $S_{(f)}\xrightarrow\sim S/(f-1)S$.
映射 $\mathfrak p\mapsto\mathfrak pS_{f,\bullet}\cap S_{(f)}$ 是从 $D(f)$ 到 $\operatorname{Spec}(S_{(f)})$ 的同胚，可借此将 $\operatorname{Spec}(S_{(f)})$ 的结构层搬到 $D(f)$ 上. 若 $f'$ 的次数为 $d'$，则 $D(f)$ 的开集 $D(ff')$ 对应于 $\operatorname{Spec}(S_{(f)})$ 中与元素 $f'^df^{-d'}$ 相伴的主开集. 由于 $S_{(ff')}$ 可等同于 $S_{(f)}$ 关于此元素的局部化，刚才在各个 $D(f)$ 上定义的层粘合成 $\operatorname{Proj}(S_\bullet)$ 上的结构层 $\mathcal O$. 这样，$\operatorname{Proj}(S_\bullet)$ 就带有概形结构，甚至是 $S_0$-概形结构.
\emph{例：}设 $S_\bullet=A[x_0,\ldots,x_n]$. 得到的概形是 $A$ 上的 $n$ 维\newterm{ag-projective-space}{射影空间}{projective space}，记为 $\mathbb P_A^n$. 开集 $D(x_i)$，$i=0,\ldots,n-1$，覆盖 $\mathbb P_A^n$，且有
\begin{align*}
D(x_i)=\operatorname{Spec}\bigl(A[x_0/x_i,\ldots,x_n/x_i]\bigr)\simeq\mathbb A_A^n.
\end{align*}
更一般地，若 $M$ 是 $A$-模，令 $S_\bullet=\operatorname{Sym}_A(M)$ 为它的\newterm{ag-symmetric-algebra}{对称代数}{symmetric algebra}. 回忆，这是 $M$ 的张量代数的``交换化''. 特别地，$\operatorname{Sym}_A(M)_n=(M^{\otimes n})_{\mathfrak S_n}$（对称群作用下的余不变量）. 此外，包含映射 $M=\operatorname{Sym}_A(M)_1\hookrightarrow\operatorname{Sym}_A(M)$ 在交换 $A$-代数范畴上表出函子 $B\mapsto\operatorname{Hom}_{A\text{-}\mathrm{Mod}}(M,B)$. 记 $\mathbb P_A(M):=\operatorname{Proj}(S_\bullet)$. 若 $(m_i)_{i\in I}$ 是 $A$-模 $M$ 的一族生成元，则 $D(m_i)$ 覆盖 $\mathbb P_A(M)$.
注意，$\operatorname{Proj}$ 的构造对分次环 $S_\bullet$ 只有部分函子性. 事实上，分次态射 $S_\bullet\xrightarrow\varphi R_\bullet$ 诱导概形态射
\begin{align*}
U_\varphi\subset\operatorname{Proj}(R_\bullet)\xrightarrow{\operatorname{Proj}(\varphi)}\operatorname{Proj}(S_\bullet),
\end{align*}
其中 $U_\varphi$ 是 $\operatorname{Proj}(R_\bullet)$ 中由不包含 $\varphi(S_+)$ 的 $\mathfrak p$ 组成的开集. 一个重要特例是 $\varphi$ 满射的情形. 此时 $U_\varphi=\operatorname{Proj}(R_\bullet)$，且上面的态射是闭浸入. 特别地，$S_\bullet$ 的每个齐次理想都定义 $\operatorname{Proj}(S_\bullet)$ 的一个闭子概形.
\emph{注：}两个理想可能定义同一个闭子概形. 例如 $\mathbb P_A^1$ 中的 $(x_0)$ 和 $(x_0^2,x_0x_1)$. 同样，也可能 $U_\varphi=\operatorname{Proj}(R_\bullet)$ 且 $\operatorname{Proj}(\varphi)$ 是同构，而 $\varphi$ 不是同构. 当 $\varphi$ 对 $n\geqslant n_0$ 诱导同构 $S_n\xrightarrow\sim R_n$ 时，便是如此. 下面是另一个重要例子.
\emph{例：}设 $S_\bullet^{(d)}$ 是由 $S_n^{(d)}:=S_{nd}$ 定义的分次环. 由构造，显然有 $\operatorname{Proj}(S_\bullet^{(d)})=\operatorname{Proj}(S_\bullet^{\langle d\rangle})$，其中分次环 $S_\bullet^{\langle d\rangle}$ 定义为：$d$ 整除 $n$ 时 $S_n^{\langle d\rangle}=S_n$，否则 $S_n^{\langle d\rangle}=0$. 可以证明，分次包含映射 $S_\bullet^{\langle d\rangle}\hookrightarrow S_\bullet$ 诱导同构 $\operatorname{Proj}(S_\bullet)\xrightarrow\sim\operatorname{Proj}(S_\bullet^{\langle d\rangle})$. 事实上，可以找到一族次数被 $d$ 整除的元素 $(f_i)_{i\in I}$，使 $D(f_i)$，$i\in I$，覆盖 $\operatorname{Proj}(S_\bullet)$. 这些 $D(f_i)$ 也覆盖 $\operatorname{Proj}(S_\bullet^{\langle d\rangle})$，且由构造有 $S_{(f_i)}=S_{(f_i)}^{\langle d\rangle}$.
\emph{应用：}例如，这使我们能定义嵌入 $\mathbb P_A^n\hookrightarrow\mathbb P_A^N$，其中 $N=\binom{n+d}{d}-1=\dim_A(A[x_0,\ldots,x_n]_d)$：将 $\mathbb P^N$ 的坐标送到 $\mathbb P^n$ 的 $d$ 次单项式.
设 $M=M_\bullet$ 是分次 $S_\bullet$-模. 对 $S_+$ 中的齐次元素 $f$，局部化 $M_{f,\bullet}$ 中次数 $0$ 的齐次元素所成的模 $M_{(f)}=M_{f,0}$ 在 $D(f)$ 上定义一个模层. 各个 $D(f)$ 上的这些层粘合成 $\operatorname{Proj}(S)$ 上的拟凝聚层 $\widetilde M$. 对整数 $n$，记 $M(n)=M(n)_\bullet$ 为将 $M_\bullet$ 的分次平移 $n$ 所得的分次模，即 $M(n)_i=M_{i+n}$. 通常记 $\mathcal O(n):=\widetilde{S(n)}$. 注意，对每个 $n$ 都有自然态射
\begin{align*}
M_n\longrightarrow\Gamma(\operatorname{Proj}(S),\widetilde{M(n)}).
\end{align*}
事实上，对每个齐次 $f\in S_+$，有 $\Gamma(D(f),\widetilde{M(n)})=M(n)_{(f)}=M(n)_{f,0}=M_{f,n}$，上面的态射由自然态射 $M_n\longrightarrow M_{f,n}$ 粘合而成.
反过来，给定 $X=\operatorname{Proj}(S_\bullet)$ 上的拟凝聚层 $\mathcal F$，为它对应分次代数 $\Gamma_\bullet(\mathcal O)$ 上的分次模 $\Gamma_\bullet(\mathcal F):=\bigoplus_{n\in\mathbb Z}\Gamma(X,\mathcal F\otimes_{\mathcal O}\mathcal O(n))$. 上面构造的分次态射 $S_\bullet\longrightarrow\Gamma_\bullet(\mathcal O)$（取 $M=S$）使每个 $\Gamma_\bullet(\mathcal F)$ 成为分次 $S_\bullet$-模. 注意，对次数为 $d$ 的齐次元素 $f$，有态射
\begin{align*}
\Gamma(D(f),\widetilde{\Gamma_\bullet(\mathcal F)})
&=\varinjlim\bigl(\Gamma(X,\mathcal F)\xrightarrow{\times f}\Gamma(X,\mathcal F\otimes\mathcal O(d))\xrightarrow{\times f}\cdots\bigr)\\
&\longrightarrow\varinjlim\bigl(\Gamma(D(f),\mathcal F)\xrightarrow{\times f}\Gamma(D(f),\mathcal F\otimes\mathcal O(d))\xrightarrow{\times f}\cdots\bigr)\\
&=\Gamma(D(f),\mathcal F)\otimes\varinjlim\bigl(S_{f,0}\xrightarrow{\times f}S_{f,d}\xrightarrow{\times f}\cdots\bigr)\\
&=\Gamma(D(f),\mathcal F).
\end{align*}
这些态射粘合成层态射 $\widetilde{\Gamma_\bullet(\mathcal F)}\longrightarrow\mathcal F$，它对 $\mathcal F$ 具有函子性.
要使这些构造具有有意义的性质，必须假设 $S_\bullet$ 作为 $S_0$-代数由次数 $1$ 的元素生成. 关键在于，此时次数 $1$ 的 $f$ 所给出的 $D(f)$ 覆盖 $\operatorname{Proj}(S_\bullet)$.
\end{numpara}
% 原PDF p17--p18上部.
\begin{proposition}\label{item:ag-2-1-2}
假设 $S_\bullet$ 作为 $S_0$-代数由 $S_1$ 生成.
\begin{enumerate}[label=\roman*)]
\item 对每个 $n$，$\operatorname{Proj}(S_\bullet)$ 上的层 $\mathcal O(n)$ 可逆. 此外，对每个分次 $S_\bullet$-模 $M_\bullet$，有典范同构 $\widetilde{M(n)_\bullet}\xrightarrow\sim\widetilde{M_\bullet}\otimes_{\mathcal O}\mathcal O(n)$.
\item 函子 $\mathcal F\mapsto\Gamma_\bullet(\mathcal F)$ 是函子 $M_\bullet\mapsto\widetilde{M_\bullet}$ 的右伴随.
\item 若再假设 $S_1$ 在 $S_0$ 上有限生成\footnote{特别地，当 $S_\bullet$ 为 Noether 环时，此条件成立.}，则对每个 $\mathcal F$，伴随箭头 $\widetilde{\Gamma_\bullet(\mathcal F)}\longrightarrow\mathcal F$ 都是同构.
\item 当 $S_\bullet=\operatorname{Sym}_{S_0}^\bullet(S_1)$ 且 $S_1$ 自由、秩为 $n+1$ 时（从而 $S_\bullet\simeq S_0[X_0,X_1,\ldots,X_n]$），伴随箭头 $S_\bullet\longrightarrow\Gamma_\bullet(\mathcal O)$ 是同构.
\end{enumerate}
\end{proposition}
\begin{proof}
\textup{(i)} 对次数 $1$ 的元素 $f$，乘以 $f^n$ 诱导同构 $\Gamma(D(f),\mathcal O)=S_{f,0}\xrightarrow\sim\Gamma(D(f),\mathcal O(n))=S_{f,n}$. 因此 $\mathcal O(n)|_{D(f)}$ 自由、秩为 $1$. 同样，此乘法诱导 $M_{f,0}\xrightarrow\sim M_{f,n}$，因而诱导 $S_{f,n}\otimes_{S_{f,0}}M_{f,0}\xrightarrow\sim M_{f,n}$，它给出所述同构在 $D(f)$ 上的限制. 再验证这些同构可以粘合.
\textup{(ii)} 留作习题. 参见 EGA II 2.6.5.
\textup{(iii)} 须证明上面定义的两个归纳极限之间的箭头是双射（还可假设 $f$ 的次数为 $1$）. 为此，只需对 $X=\operatorname{Proj}(S_\bullet)$ 上的凝聚层 $\mathcal G$ 应用下面两个性质：
\begin{enumerate}[label=\roman*)]
\item 对每个截面 $s\in\operatorname{Ker}(\Gamma(X,\mathcal G)\longrightarrow\Gamma(D(f),\mathcal G))$，存在 $n$ 使 $s\otimes f^n=0$ 于 $\Gamma(X,\mathcal G\otimes\mathcal O(n))$ 中.
\item 对每个截面 $s\in\Gamma(D(f),\mathcal G)$，存在整数 $n$，使 $s\otimes f^n$ 可延拓为 $\Gamma(X,\mathcal G\otimes\mathcal O(n))$ 中的截面.
\end{enumerate}
这些性质留作习题（参见 Hartshorne，引理 II.5.14）. $S_1$ 在 $S_0$ 上有限生成的假设，用来将 $X$ 用有限多个形如 $D(f)$、$f$ 次数为 $1$ 的仿射开集覆盖.
\textup{(iv)} 利用正合列计算 $\Gamma_\bullet(\mathcal O)$：
\begin{align*}
\begin{tikzcd}[ampersand replacement=\&,column sep=small,row sep=small]
0\arrow[r]\&\Gamma_\bullet(\mathbb P_{S_0}^n,\mathcal O)\arrow[r]\&\prod_{i=0}^n\Gamma_\bullet(D(X_i),\mathcal O)\arrow[r]\arrow[d,equal]\&\prod_{i,j}\Gamma_\bullet(D(X_i)\cap D(X_j),\mathcal O)\arrow[d,equal]\\
\&\&\prod_{i=0}^nS_{X_i,\bullet}\arrow[r]\&\prod_{i,j}S_{X_iX_j,\bullet}
\end{tikzcd}
\end{align*}
它表明 $\Gamma_\bullet(\mathcal O)$ 是各局部化 $S_{X_i}$ 在局部化 $S_{X_0X_1\cdots X_n}$ 中的交. 此局部化中的元素唯一地写成 $\prod_{i=0}^nX_i^{a_i}P(X_0,\ldots,X_n)$，其中 $P$ 不被 $X_0,\ldots,X_n$ 中的任何一个整除. 它属于 $S_{X_i}$ 当且仅当 $a_j$，$j\ne i$，均非负，故它属于各个 $S_{X_i}$ 的交当且仅当所有 $a_i$ 均非负，也即当且仅当它属于 $S$（参见 Hartshorne，命题 5.13）.
\end{proof}
\begin{remark}\label{item:ag-2-1-3}
\textup{(i)} 命题的 ii) 表明 $\Gamma_\bullet$ 是全忠实函子，且 $M\mapsto\widetilde M$ 本质满射. 然而，前者并非本质满射，后者也不忠实. 例如，对任意整数 $n_0$，若 $M'_\bullet$ 表示分次模 $\bigoplus_{n\geqslant n_0}M_n$，则包含映射 $M'_\bullet\hookrightarrow M_\bullet$ 诱导相伴层的同构 $\widetilde{M'}\xrightarrow\sim\widetilde M$. 这可在每个次数为 $1$ 的 $f$ 所对应的 $D(f)$ 上验证：$M'_{(f)}$ 是某个归纳系统的归纳极限，而此系统在以 $M_{(f)}$ 为归纳极限的系统中共尾.
\textup{(ii)} 若 $A$ 是环，由 iii)、iv) 以及 $\Gamma_\bullet$ 左正合可知，射影空间 $\mathbb P_A^n:=\operatorname{Proj}(\operatorname{Sym}^\bullet(A^{n+1}))=\operatorname{Proj}(A[X_0,X_1,\ldots,X_n])$ 的每个闭子概形都由齐次理想 $I$ 定义，例如可取 $I=\Gamma_\bullet(\mathcal I)$，其中 $\mathcal I$ 是定义该闭子概形的理想层.
\textup{(iii)} 还应注意，在命题的假设下，有分次 $S_\bullet$-模的满态射 $S_\bullet\otimes_{S_0}S_1\twoheadrightarrow S(1)_{\bullet\geqslant0}$，因此它在 $\operatorname{Proj}(S_\bullet)$ 上给出层的满态射
\begin{align}
\mathcal O\otimes_{S_0}S_1\twoheadrightarrow\mathcal O(1).\qquad\diamond\tag{2.1.3.1}\label{eq:ag-2-1-3-1}
\end{align}
\renewcommand{\envenddiamond}{}% marker is on the final display
\end{remark}
% 原PDF p18--p19.
\begin{numpara}[整体化]\label{item:ag-2-1-4}
设 $S$ 是概形，$\mathcal S_\bullet=\mathcal O_S\oplus\mathcal S_1\oplus\cdots$ 是分次 $\mathcal O_S$-代数的拟凝聚层. 用仿射开集覆盖 $S$，通过粘合定义 $S$ 上的概形 $\operatorname{Proj}(\mathcal S_\bullet)$. 更确切地说：
\begin{enumerate}[label=\roman*)]
\item 先验证，若 $S=\operatorname{Spec}(A)$ 且 $U=\operatorname{Spec}(A_f)$，则有典范的 $U$-同构 $\operatorname{Proj}(S_\bullet\otimes_AA_f)\xrightarrow\sim\operatorname{Proj}(S_\bullet)\times_XU$.
\item 当 $S$ 分离时，选择仿射覆盖 $S=\bigcup_{U\in\mathcal U}U$，利用 i) 粘合各个 $\operatorname{Proj}(\mathcal S_\bullet(U))$. 仍利用 i)，验证所得概形与覆盖的选择无关.
\item 对一般的 $S$，选择由分离开集（例如仿射开集）组成的覆盖，利用 ii) 粘合各个 $\operatorname{Proj}(\mathcal S_\bullet|_U)$.
\end{enumerate}
\emph{例：}设 $\mathcal F$ 是 $\mathcal O_S$-模. 记 $\operatorname{Sym}_{\mathcal O_S}(\mathcal F)$ 为预层 $U\mapsto\operatorname{Sym}_{\mathcal O_S(U)}(\mathcal F(U))$ 的相伴层. 这是一个分次 $\mathcal O_S$-代数，带有 $\mathcal O_S$-模态射 $\mathcal F\hookrightarrow\operatorname{Sym}_{\mathcal O_S}(\mathcal F)$，它在 $\mathcal O_S$-代数范畴上表出函子 $\operatorname{Hom}_{\mathcal O_S\text{-}\mathrm{mod}}(\mathcal F,-)$. 若 $\mathcal F$ 拟凝聚，则 $\operatorname{Sym}_{\mathcal O_S}(\mathcal F)$ 也拟凝聚. 事实上，只需在仿射情形 $S=\operatorname{Spec}(A)$ 下验证. 此时令 $M=\Gamma(S,\mathcal F)$. 对每个 $\mathcal O_S$-代数 $\mathcal A$，有
\begin{align*}
\operatorname{Hom}_{\mathcal O_S\text{-}\mathrm{mod}}(\mathcal F,\mathcal A)
&\xrightarrow\sim\operatorname{Hom}_{A\text{-}\mathrm{mod}}(M,\Gamma(S,\mathcal A))\\
&\simeq\operatorname{Hom}_{A\text{-}\mathrm{alg}}(\operatorname{Sym}_A(M),\Gamma(S,\mathcal A)).
\end{align*}
这表明，与 $A$-代数 $\operatorname{Sym}_A(M)$ 相伴的拟凝聚 $\mathcal O_S$-代数表出与 $\operatorname{Sym}_{\mathcal O_S}(\mathcal F)$ 相同的函子，所以与它同构. 还应注意以下性质：
\begin{itemize}
\item 对每个 $s\in S$，有 $\operatorname{Sym}_{\mathcal O_S}(\mathcal F)_s=\operatorname{Sym}_{\mathcal O_{S,s}}(\mathcal F_s)$.
\item 对每个 $f:X\longrightarrow S$，有 $f^*\operatorname{Sym}_{\mathcal O_S}(\mathcal F)\xrightarrow\sim\operatorname{Sym}_{\mathcal O_X}(f^*\mathcal F)$.
\end{itemize}
它们都由换基同构 $\operatorname{Sym}_A(M)\otimes_AA'\simeq\operatorname{Sym}_{A'}(M\otimes_AA')$ 得出；此同构对每个环态射 $A\longrightarrow A'$ 成立.
于是记
\begin{align*}
\mathbb P(\mathcal F):=\operatorname{Proj}(\operatorname{Sym}_{\mathcal O_S}(\mathcal F)),\qquad
\pi:\mathbb P(\mathcal F)\longrightarrow S
\end{align*}
为典范态射. 当 $\mathcal F=\mathcal O_S^{n+1}$ 时，有 $\mathbb P(\mathcal F)\simeq\mathbb P_S^n=\mathbb P^n\times S$. 当 $\mathcal F$ 在 $S$ 上局部自由、秩为 $n+1$ 时，$S$-概形 $\mathbb P(\mathcal F)$ 局部地（在 $S$ 上）同构于 $n$ 维射影空间. 这是更一般的换基同构 $\mathbb P(f^*\mathcal F)\xrightarrow\sim\mathbb P(\mathcal F)\times_SX$ 的推论，此同构对每个 $f:X\longrightarrow S$ 成立. 特别地，若 $\mathcal F$ 是可逆层，则态射 $\pi_{\mathcal F}:\mathbb P(\mathcal F)\longrightarrow S$ 是同构. 对任意 $\mathcal F$，满态射~\eqref{eq:ag-2-1-3-1}粘合成 $\mathcal O_{\mathbb P(\mathcal F)}$-模层的满态射
\begin{align}
\pi^*\mathcal F\twoheadrightarrow\mathcal O_{\mathbb P(\mathcal F)}(1).\tag{2.1.4.1}\label{eq:ag-2-1-4-1}
\end{align}
应当牢记 $\mathbb P(\mathcal F)$ 的模空间解释：
\end{numpara}
\begin{proposition}\label{item:ag-2-1-5}
设 $S$ 是概形，$\mathcal F$ 是拟凝聚 $\mathcal O_S$-模. 对每个 $S$-概形 $X\xrightarrow f S$，映射
\begin{align*}
\operatorname{Hom}_S(X,\mathbb P(\mathcal F))
&\longrightarrow\{\text{商 }f^*\mathcal F\twoheadrightarrow\mathcal L,\ \mathcal L\text{ 为可逆 }\mathcal O_X\text{-模}\},\\
g:X\longrightarrow\mathbb P(\mathcal F)
&\longmapsto\bigl(f^*\mathcal F\twoheadrightarrow g^*\mathcal O_{\mathbb P(\mathcal F)}(1)\bigr)
\end{align*}
是双射.
\end{proposition}
\emph{说明：}在右边，若存在同构 $\varphi:\mathcal L\xrightarrow\sim\mathcal L'$ 使得 $p'=\varphi\circ p$，则将两个商 $p:f^*\mathcal F\twoheadrightarrow\mathcal L$ 和 $p':f^*\mathcal F\twoheadrightarrow\mathcal L'$ 等同.
\begin{proof}
定义逆双射. 设 $f^*\mathcal F\twoheadrightarrow\mathcal L$ 是商，其中 $\mathcal L$ 是 $X$ 上的可逆层. 由此得到满的分次态射 $\operatorname{Sym}_{\mathcal O_X}^\bullet(f^*\mathcal F)\xrightarrow\varphi\operatorname{Sym}_{\mathcal O_X}^\bullet(\mathcal L)$，于是由函子性得到 $X$-态射
\begin{align*}
U_\varphi\subset\mathbb P(\mathcal L)=X\longrightarrow\mathbb P(f^*\mathcal F)=X\times_S\mathbb P(\mathcal F).
\end{align*}
由于 $\varphi$ 满射，故 $U_\varphi=\mathbb P(\mathcal L)$；再作投影，就得到所求态射 $X\longrightarrow\mathbb P(\mathcal F)$. 容易验证，这给出了命题中双射的逆.
\end{proof}
% 原PDF p19下部--p20.
\begin{remark}\label{item:ag-2-1-6}
当 $\mathcal F$ 凝聚时，利用赋值判别法，上述模空间解释立即表明 $\mathbb P(\mathcal F)$ 在 $S$ 上紧合. 事实上，设 $R$ 是分式域为 $K$ 的赋值环，$\operatorname{Spec}(R)\xrightarrow x S$ 是一个 $R$-点. ``纤维'' $\mathcal F_R:=x^*\mathcal F$ 是有限生成 $R$-模，纤维 $\mathcal F_K=K\otimes \mathcal F_R$ 是有限维 $K$-向量空间. $x|_{\operatorname{Spec}(K)}$ 上方的 $\mathbb P(\mathcal F)$ 的一个 $K$-点，就是满态射 $\mathcal F_K\twoheadrightarrow K$. 而复合映射 $\mathcal F_R\longrightarrow \mathcal F_K\longrightarrow K$ 的像是 $K$ 的有限生成 $R$-子模，因而同构于 $R$. 所以此复合映射给出 $x$ 上方的 $\mathbb P(\mathcal F)$ 的一个 $R$-点.
稍后我们将看到如何利用这种模空间解释来检验态射 $\pi$ 的光滑性.\unskip\nobreak\hfill\ensuremath{\diamond}
\renewcommand{\envenddiamond}{}% marker is on the final text line
\end{remark}
回忆，浸入是若干闭浸入与开浸入的复合.
\begin{definition}\label{item:ag-2-1-7}
设 $f:X\longrightarrow S$ 是概形态射.
\begin{enumerate}[label=\roman*)]
\item 若 $f$ 可经由到某个射影空间 $\mathbb P_S^n:=\mathbb P_{\mathbb Z}^n\times S$ 的浸入（分别为闭浸入）分解，则称 $f$ 是\newterm{ag-quasi-projective-morphism}{拟射影态射}{quasi-projective morphism}（分别为\newterm{ag-projective-morphism}{射影态射}{projective morphism}），也称 $X$ 在 $S$ 上拟射影（分别为射影）.
\item 对任意 $S$-浸入 $i:X\hookrightarrow\mathbb P_S^r$，称形如 $i^*(\mathcal O_{\mathbb P^r}(1))$ 的层为 $X$ 上相对于 $S$ 的\newterm{ag-relatively-very-ample-sheaf}{极丰沛层}{relatively very ample sheaf}.
\end{enumerate}
\end{definition}
若 $S=\operatorname{Spec}(A)$，则每个形如 $\operatorname{Proj}(S_\bullet)$ 的 $S$-概形都在 $S$ 上射影，其中 $S_\bullet$ 是满足 $S_0=A$、且由有限生成 $A$-模 $S_1$ 生成的分次环. 此时，$\operatorname{Proj}(S_\bullet)$ 上的层 $\mathcal O(1)=\widetilde{S(1)_\bullet}$ 相对于 $A$ 极丰沛. 反过来，$S$ 上的每个射影概形都有此形式，因为此时可找到满的分次态射 $A[T_0,\ldots,T_n]\twoheadrightarrow S_\bullet$.
然而，若底概形 $S$ 非仿射，而 $\mathcal S_\bullet$ 是分次 $\mathcal O_S$-代数，则相伴的 $\operatorname{Proj}(\mathcal S_\bullet)$ 上的层 $\mathcal O(1)$ 未必相对于 $S$ 极丰沛. 例如，取 $\mathcal S_\bullet=\operatorname{Sym}^\bullet(\mathcal L^2)$，其中 $\mathcal L$ 是 $S$ 上没有非零整体截面的层. 到第二个因子的投影 $\mathcal L^2\twoheadrightarrow\mathcal L$ 定义截面 $s:S\hookrightarrow\mathbb P(\mathcal L^2)$，按定义有 $s^*\mathcal O_{\mathbb P(\mathcal L^2)}(1)\xrightarrow\sim\mathcal L$. 但这表明，$\mathcal O_{\mathbb P(\mathcal L^2)}(1)$ 不可能由其整体截面生成，所以不可能相对于 $S$ 极丰沛.
\begin{remark}\label{item:ag-2-1-8}
由前面的命题，态射 $X\xrightarrow g\mathbb P_S^n=\operatorname{Proj}(\mathcal O_S[T_0,\ldots,T_n])$ 由 $X$ 上的一个可逆层 $\mathcal L$ 和生成 $\mathcal L$ 的 $n+1$ 个截面 $s_0,\ldots,s_n$ 给出（即相伴态射 $\mathcal O_X^{n+1}\longrightarrow\mathcal L$ 满射）. 从 $g$ 出发，有 $\mathcal L=g^*\mathcal O(1)$ 及 $s_i=g^*(T_i)$. 反过来，从 $\mathcal L$ 和 $s_i$ 出发，设 $X_{s_i}$ 是 $s_i$ 不为零的开集. 则有平凡化 $s_i:\mathcal O_{X_{s_i}}\xrightarrow\sim\mathcal L|_{X_{s_i}}$，由此得到 $n$ 个截面 $s_j/s_i\in\Gamma(X_{s_i},\mathcal O_X)$，它们定义态射
\begin{align*}
X_{s_i}\xrightarrow{g_i}D(T_i)=\mathbb A^n\subset\mathbb P^n.
\end{align*}
态射 $g$ 由各个 $g_i$ 粘合而成. 注意 $X_{s_i}=g^{-1}(D(T_i))$. 再假设 $S$ 仿射，等于 $\operatorname{Spec}(A)$. 则 $g$ 是闭浸入，当且仅当：
\begin{enumerate}[label=\roman*)]
\item 各个 $X_{s_i}$ 都仿射.
\item 截面 $s_j/s_i$ 作为 $A$-代数的生成元生成 $\Gamma(X_{s_i},\mathcal O_X)$.\unskip\nobreak\hfill\ensuremath{\diamond}
\end{enumerate}
\renewcommand{\envenddiamond}{}% marker is on the last list item
\end{remark}
% 原PDF p21.
\begin{definition}\label{item:ag-2-1-9}
设 $X$ 是概形. 若对 $X$ 上的每个凝聚层 $\mathcal F$，都存在 $n_0$，使得对每个 $n\geqslant n_0$，层 $\mathcal F\otimes\mathcal L^{\otimes n}$ 由有限多个整体截面生成（即存在满态射 $\mathcal O_X^N\twoheadrightarrow\mathcal F\otimes\mathcal L^{\otimes n}$），则称可逆层 $\mathcal L$ 为\newterm{ag-ample-sheaf}{丰沛层}{ample sheaf}.
\end{definition}
\begin{theorem}\label{item:ag-2-1-10}
设 $X$ 是 Noether 环 $A$ 上的有限型概形，$\mathcal L$ 是 $X$ 上的可逆层. 则 $\mathcal L$ 丰沛，当且仅当存在 $n$ 使 $\mathcal L^{\otimes n}$ 相对于 $A$ 极丰沛.
\end{theorem}
由此可知，概形 $X$ 在 $A$ 上拟射影，当且仅当它具有丰沛层.
\begin{proof}
假设 $\mathcal L$ 丰沛. 先证明，存在整数 $n$ 和截面 $s_1,\ldots,s_k\in\Gamma(X,\mathcal L^{\otimes n})$，使开集 $X_{s_i}$ 都仿射并覆盖 $X$（所以 $s_i$ 生成 $\mathcal L^{\otimes n}$）. 事实上，设 $x\in X$，取包含 $x$ 的仿射开集 $U$，使 $\mathcal L$ 在 $U$ 上自由. 设 $\mathcal I$ 是其闭补集的理想层. 因为 $\mathcal L$ 丰沛，存在 $n$ 使 $\mathcal I\otimes\mathcal L^{\otimes n}$ 由整体截面生成. 特别地，存在截面 $s\in\Gamma(X,\mathcal I\otimes\mathcal L^{\otimes n})\subset\Gamma(X,\mathcal L^{\otimes n})$，使 $x\in X_s\subset U$. 由于 $\mathcal L$ 在 $U$ 上自由，存在函数 $f\in\Gamma(U,\mathcal O_U)$（即 $s|_U$ 在平凡化 $\mathcal L|_U\xrightarrow\sim\mathcal O_U$ 下的像），使 $X_s=U_f$，故 $X_s$ 仿射. 由 $X$ 的拟紧性，可用有限多个这种开集覆盖 $X$. 注意 $X_s=X_{s^n}$，因此可以为所有这些截面找到共同的 $n$.
现在对每个 $i$，固定 $A$-代数 $\Gamma(X_{s_i},\mathcal O_X)$ 的有限多个生成元 $f_{ij}$. 重用命题~\ref{item:ag-2-1-2}证明中的论证，存在整数 $m$，使每个 $f_{ij}\otimes s_i^m$ 延拓为截面 $t_{ij}\in\Gamma(X,\mathcal L^{\otimes nm})$. 考虑由可逆层 $\mathcal L^{\otimes nm}$ 以及所有的截面 $s_i^m$、$t_{ij}$ 所确定的态射 $X\xrightarrow g\mathbb P^N$. 记 $\mathbb P^N$ 上相应的坐标为 $S_i$ 和 $T_{ij}$. 则由构造：
\begin{itemize}
\item $g$ 经由 $\mathbb P_A^N$ 的开集 $\bigcup_iD(S_i)$ 分解.
\item 限制 $g|_{X_{s_i}}:X_{s_i}\longrightarrow D(S_i)$ 是闭浸入，因为态射 $\Gamma(U_i,\mathcal O_{\mathbb P_A^N})=A[S_i,T_{i,j}]\longrightarrow\Gamma(X_{s_i},\mathcal O_X)$ 满射.
\end{itemize}
因此 $g$ 是 $X$ 到 $\mathbb P_A^N$ 的浸入，故 $\mathcal L^{\otimes nm}$ 相对于 $A$ 极丰沛.
现在假设 $\mathcal L$ 相对于 $A$ 极丰沛，并选取定义到 $\mathbb P_A^n$ 的浸入的截面 $s_0,\ldots,s_n$. 先假设此浸入为闭浸入（即 $X$ 紧合）. 则 $X_{s_i}$ 仿射（参见上面的注）. 取 $\Gamma(X_{s_i},\mathcal F)$ 作为 $\Gamma(X_{s_i},\mathcal O_X)$-模的生成元 $m_{ij}$. 仍用命题~\ref{item:ag-2-1-2}证明中的论证，存在整数 $n$，使截面 $m_{ij}s_i^n$ 延拓为 $\Gamma(X,\mathcal F\otimes\mathcal L^{\otimes n})$ 中的整体截面. 由构造，这些截面生成 $\mathcal F\otimes\mathcal L^{\otimes n}$. 回到一般情形（$X$ 未必紧合），记 $\overline X$ 为 $X$ 在 $\mathbb P_A^n$ 中的闭包，$\overline{\mathcal L}$ 为 $\mathcal O_{\mathbb P_A^n}(1)$ 在 $\overline X$ 上的限制. 则 $X$ 是 $\overline X$ 的开集；断言由下述事实得出：$X$ 上的每个凝聚层 $\mathcal F$ 都可延拓为 $\overline X$ 上的凝聚层，参见 Hartshorne，习题 II.5.15.
\end{proof}
% 原PDF p22--p23上部（止于2.2之前）.
\begin{theorem}\label{item:ag-2-1-11}
设 $A$ 是域上（或更一般地，Nagata 环上）的有限生成代数，$\mathcal F$ 是 $A$ 上射影概形 $X$ 上的凝聚层. 则 $\Gamma(X,\mathcal F)$ 作为 $A$-模有限生成.
\end{theorem}
当 $A$ 为 Noether 环时，结论仍成立，但通常的证明（归于 Serre）用到拟凝聚层的上同调，参见 Hartshorne III.5.2 或 EGA III.2.1（即使只关心 $H^0$，也要对 $i$ 作下降归纳，建立对所有 $H^i$ 都成立的结果）.
\newterm{ag-nagata-ring}{Nagata 环}{Nagata ring}是``普遍日本''的 Noether 整环，即 $A$ 上的每个有限生成整代数都是``日本''的. 若对分式域 $Q(A)$ 的每个有限扩张 $K$，$A$ 在 $K$ 中的整闭包在 $A$ 上有限生成（从而因其为整的，故在 $A$ 上有限），则称整环 $A$ 为\newterm{ag-japanese-ring}{日本环}{Japanese ring}. 每个域都是 Nagata 环，环 $\mathbb Z$ 也是；更一般地，每个分式域为完美域的 Noether 整闭环都是 Nagata 环.
\begin{proof}
这里只给出证明的步骤. 参见 Hartshorne，定理 5.19. 由假设，$X=\operatorname{Proj}(S_\bullet)$，其中 $S_0=A$，且 $S_\bullet$ 由有限生成 $A$-模 $S_1$ 生成.
\begin{enumerate}[label=\arabic*-]
\item 证明 $\mathcal F=\widetilde M$，其中 $M=M_\bullet$ 是有限生成分次 $S_\bullet$-模. 事实上，只需在 $\Gamma_\bullet(\mathcal F)$ 中取由某个 $\mathcal F\otimes\mathcal O(n)$ 的有限多个生成截面所生成的分次子模.
\item 归结到 $S$ 为整环且 $M_\bullet$ 形如 $S_\bullet(n)$ 的情形. 这里利用 $M$ 有一个滤过，其分次商模形如 $(S_\bullet/\mathfrak p)(n)$.
\item 假设 $S$ 为整环，证明 $S':=\Gamma_\bullet(\mathcal O)$ 在 $S$ 上整. 事实上，固定生成元 $f_0,\ldots,f_n\in S_1$. 用 $D(f_i)$ 覆盖 $X$，得到正合列
\begin{align*}
0\longrightarrow\Gamma_\bullet(\mathcal O)\longrightarrow\bigoplus_{i=0}^nS_{f_i,\bullet}\longrightarrow\bigoplus_{i,j=0}^nS_{f_if_j,\bullet}.
\end{align*}
它表明 $\Gamma_\bullet(\mathcal O)$ 可等同于各 $S_{f_i}$ 在 $S_{f_1\cdots f_0}$ 中的交. 只需验证 $\bigcap_iS_{f_i}$ 在 $S$ 上整. 设 $f'$ 属于此交，选择整数 $m$ 使每个 $i$ 都有 $f_i^mf'\in S$. 从而 $S_{\geqslant nm}f'\subset S$，甚至 $S_{\geqslant nm}f'\subset S_{\geqslant nm}$，因为 $f'$ 的次数为正. 由归纳法，对每个 $d\geqslant0$，有 $S_{\geqslant nm}f'^d\subset S_{\geqslant nm}$. 因而
\begin{align*}
Sf'^d\subset f_0^{-nm}S_{\geqslant nm}f'^d\subset f_0^{-nm}S_{\geqslant nm}\subset f_0^{-nm}S,
\end{align*}
最终有 $S[f']\subset f_0^{-nm}S$. 因为 $S$ 为 Noether 环，所以 $S[f']$ 在 $S$ 上有限，因而 $f'$ 是整的.
\item 利用 $S$ 为日本环（即其整闭包作为 $S$-模有限生成）得出结论.\qedhere
\end{enumerate}
\end{proof}
\emph{注：}沿用证明中的记号，可知 $S'=\Gamma_\bullet(\mathcal O)$ 包含在 $S$ 在其关于 $S_1$ 中元素的局部化中的整闭包内. 由我们的假设，$S'$ 在 $S$ 上有限. $S'\subset S_{S_1}$ 意味着，对每个元素 $f'\in S'$，存在 $n$ 使 $f'S_{\geqslant n}\subset S_{\geqslant n}$. 由 $S'$ 在 $S$ 上有限可知，存在整数 $N$，使每个 $f'\in S'$ 都有 $f'S_{\geqslant N}=S_{\geqslant N}$. 因此，存在整数 $M$ 使 $S'_{\geqslant M}=S_{\geqslant M}$. 换言之，代数 $S_\bullet$ 与 $\Gamma_\bullet(\mathcal O)$ 在充分高的次数上一致. 如分目~\ref{item:ag-2-1-1}中的例子（原书第16页），记 $S_\bullet^{(d)}$ 为代数 $\bigoplus_nS_{nd}$. 当 $d$ 充分大时，便有同构 $\operatorname{Proj}(S'^d)\xrightarrow\sim\operatorname{Proj}(S^{(d)})\xrightarrow\sim\operatorname{Proj}(S)$.
\emph{注：}以 Noether 概形 $S$（或 Nagata 环上的有限型概形）代替 $\operatorname{Spec}(A)$，并取射影态射 $f:X\longrightarrow S$，即可将此结果整体化. 结论是 $f_*(\mathcal F)$ 在 $S$ 上凝聚.
